\section*{अध्याय \ref{ch:b3:surfaces-catalysis} --- पृष्ठ, अधिशोषण और विषमांगी उत्प्रेरण}

\begin{solution}{exo:b3:surfaces-catalysis:1}
$K = \theta/((1 - \theta)p) = 0.40/(0.60 \times 2.0) = \qty{0.33}{kPa^{-1}}$; \qty{10}{kPa} पर $\theta = 3.33/4.33 = 0.77$।
\end{solution}

\begin{solution}{exo:b3:surfaces-catalysis:2}
$-20$: \omterm{def:b3:surfaces-catalysis:physi-chemi}{भौतिक अधिशोषण} (संघनन एन्थैल्पी की कोटि); $-150$: \omterm{def:b3:surfaces-catalysis:physi-chemi}{रासायनिक अधिशोषण}, जो
अकेला वियोजी हो सकता है।
\end{solution}

\begin{solution}{exo:b3:surfaces-catalysis:3}
(100): एक शीर्षस्थ, दो सेतु (चार आबंध, प्रत्येक दो परमाणुओं द्वारा साझा), एक चतुर्गुना
गर्त (चार गर्त, प्रत्येक चार परमाणुओं द्वारा साझा)। (111): एक शीर्षस्थ, तीन सेतु, दो
त्रिगुना गर्त (प्रत्येक परमाणु के चारों ओर छह त्रिभुज, प्रत्येक तीन द्वारा साझा)।
\end{solution}

\begin{solution}{exo:b3:surfaces-catalysis:4}
$\num{3.0e-6}/\num{1.5e-5} = \qty{0.20}{s^{-1}}$।
\end{solution}

\begin{solution}{exo:b3:surfaces-catalysis:5}
$\theta/(1 - \theta) = (Kp)^{1/2}$: \qty{1.0}{kPa} पर $0.111$, \qty{4.0}{kPa} पर दुगुना: $0.222$, अतः $\theta = 0.18$।
\end{solution}

\begin{solution}{exo:b3:surfaces-catalysis:6}
हर $1 + 10 + 2.5 = 13.5$: $\theta_{\ce{CO}} = 0.74$, $\theta_{\ce{O2}} = 0.19$, मुक्त स्थल 0.07।
लैंगम्यूर--हिंशेलवुड दर, जो $\theta_{\ce{CO}}\theta_{\ce{O2}}$ के समानुपाती है, CO से भरे पृष्ठ पर
ऑक्सीजन की कमी से सीमित होती है।
\end{solution}

\begin{solution}{exo:b3:surfaces-catalysis:7}
$\Delta_{\mathrm{ads}}H = -R\ln(8.0/1.0)/(1/300 - 1/350) = -8.314 \times 2.079/\num{4.76e-4} = \qty{-36}{kJ/mol}$।
\end{solution}

\begin{solution}{exo:b3:surfaces-catalysis:8}
$n_{\mathrm m} = 120/22\,414 = \qty{5.35e-3}{mol}$; क्षेत्रफल $= \num{5.35e-3} \times \num{6.022e23} \times \num{0.162e-18} =
\qty{522}{m^2.g^{-1}}$।
\end{solution}

\begin{solution}{exo:b3:surfaces-catalysis:9}
लैंगम्यूर--हिंशेलवुड: दोनों अभिकारकों का पड़ोसी स्थलों पर अधिशोषित होना आवश्यक है;
CO दृढ़ता से बँधता है और उच्च दाब पर ऑक्सीजन के लिए कोई स्थान नहीं छोड़ता।
\end{solution}

\begin{solution}{exo:b3:surfaces-catalysis:10}
निम्न आच्छादन: $100 - 60 = \qty{40}{kJ/mol}$; उच्च आच्छादन: \qty{100}{kJ/mol}। $\ln r$ का
$1/T$ के विरुद्ध आलेख निम्न ताप पर, जहाँ पृष्ठ भरा होता है, तीव्र (ढलान $-100/R$) होता है,
और उच्च ताप पर, जहाँ वह खाली होता है, कम तीव्र (ढलान $-40/R$): दोनों के बीच
मुड़ता हुआ एक वक्र।
\end{solution}

\begin{solution}{exo:b3:surfaces-catalysis:11}
$4 \times 0.10 = 0.40$ स्थल अवरुद्ध हैं: एक स्थल के लिए सक्रियता का 60\,\% शेष रहता है,
दो संलग्न मुक्त स्थलों के लिए लगभग $0.60^2 = 36\,\%$।
\end{solution}

\begin{solution}{exo:b3:surfaces-catalysis:12}
बीच वाली (\qty{-250}{kJ/mol}): पहली नाइट्रोजन को इतनी क्षीणता से बाँधती है कि \ce{N2} आसानी से
नहीं टूटता, तीसरी इतनी दृढ़ता से कि नाइट्रोजन परमाणु पृष्ठ पर टिके रहते हैं
और उसे अवरुद्ध कर देते हैं।
\end{solution}

\begin{solution}{pb:b3:surfaces-catalysis:1}
\textbf{1.} वायुमंडलीय दाब पर नाइट्रोजन लगभग \qty{77}{K} पर संघनित होती है: आपेक्षिक दाबों का पूरा परास
प्राप्त होता है, और बहुपरतें बनती हैं।
\textbf{2.} 0.05 से नीचे पृष्ठीय विषमांगता, और 0.30 से ऊपर रंध्रों में
संघनन, BET मान्यताओं का उल्लंघन करते हैं।
\textbf{3.} \num{1.278e-3}, \num{2.410e-3}, \num{3.453e-3}, \num{4.621e-3}, \num{5.659e-3}, \num{6.813e-3} \unit{g.cm^{-3}}।
\textbf{4.} $v_{\mathrm m} = 1/(0.02205 + 0.000180) = \qty{45.0}{cm^3.g^{-1}}$।
\textbf{5.} $c = 1 + 0.02205/0.000180 = 124$: पहली परत द्रव की अपेक्षा कहीं अधिक दृढ़ता से
बँधती है।
\textbf{6.} यह ढलान गुणा $x$ की तुलना में अत्यंत छोटा है, अतः इसकी आपेक्षिक त्रुटि बड़ी है;
परंतु $v_{\mathrm m}$ ढलान और अंतःखंड के योग पर निर्भर करता है, जिसमें ढलान प्रभावी है।
\textbf{7.} $45.0/22\,414 = \qty{2.01e-3}{mol.g^{-1}}$।
\textbf{8.} $\num{2.01e-3} \times \num{6.022e23} \times \num{0.162e-18} = \qty{196}{m^2.g^{-1}}$।
\textbf{9.} ऐलुमिना आधार: अनावृत प्लैटिनम (भाग III) का क्षेत्रफल लगभग \qty{0.9}{m^2.g^{-1}} है।
\textbf{10.} आण्विक चौड़ाई के रंध्र बहुत निम्न दाब पर पूरी तरह भर जाते हैं, परत
दर परत नहीं, अतः किसी \omterm{def:b3:surfaces-catalysis:coverage}{एकलपरत} की पहचान नहीं हो सकती।
\textbf{11.} $2 \times 0.205/22\,414 = \qty{1.83e-5}{mol.g^{-1}}$।
\textbf{12.} $0.0100/195.08 = \qty{5.13e-5}{mol.g^{-1}}$।
\textbf{13.} $D = 1.83/5.13 = 0.357$।
\textbf{14.} $a^3/4 = (0.3923)^3/4 = \qty{0.01509}{nm^3}$।
\textbf{15.} $\frac{\sqrt3}{4}(0.3923)^2 = \qty{0.0666}{nm^2}$, सबसे सघन फलक; (100) और (110) से 0.0770 और
\qty{0.1088}{nm^2} मिलते हैं, और तीनों स्थल घनत्वों का माध्य \qty{0.0807}{nm^2} देता है।
\textbf{16.} एक गोले में $\pi d^3/6v_{\mathrm{at}}$ परमाणु होते हैं, जिनमें से $\pi d^2/a_{\mathrm{at}}$ पृष्ठ पर हैं:
$D = 6v_{\mathrm{at}}/(a_{\mathrm{at}}d)$।
\textbf{17.} $d = 6 \times 0.01509/(0.0807 \times 0.357) = \qty{3.1}{nm}$।
\textbf{18.} प्रति पृष्ठीय Pt एक H; एक ही आकार के गोलाकार कण (परिणाम एक
पृष्ठ-भारित माध्य है); सारा प्लैटिनम अपचयित और सुलभ; आधार पर कोई हाइड्रोजन
अतिप्रवाह नहीं।
\textbf{19.} $\num{2.0e-5}/\num{1.83e-5} = \qty{1.1}{s^{-1}}$।
\textbf{20.} प्रति ग्राम प्लैटिनम $\num{2.0e-5}/0.0100 = \qty{2.0e-3}{mol.g^{-1}.s^{-1}}$।
\textbf{21.} 0.357 के स्थान पर $D = 1.12/10 = 0.112$: 3.2 का गुणक।
\textbf{22.} कि अभिक्रिया संरचना-संवेदी है: इसके \omterm{def:b3:complex-kinetics:enzyme}{सक्रिय स्थल} (किनारे,
कोने, विशेष फलक) पृष्ठ का कोई नियत अंश नहीं हैं।
\textbf{23.} यह प्रति \omterm{def:b3:complex-kinetics:enzyme}{सक्रिय स्थल} घटनाओं को गिनती है, जिससे यह प्रभाव हट जाता है कि कितनी
धातु अनावृत है।
\textbf{24.} \textbf{माध्य प्लैटिनम कण व्यास $\approx \qty{3.1}{nm}$।}
\end{solution}
